\documentclass[10pt,a4paper]{amsart}
\usepackage[margin=19mm]{geometry}
\usepackage{amsmath,amssymb,mathtools}
\usepackage[scr=rsfs,cal=euler]{mathalfa}
\usepackage{xfrac}
\usepackage[unicode,hidelinks,bookmarks=false]{hyperref}
\newcommand{\ie}{i.e.\ }
\newcommand{\cf}{cf.\ }
\newcommand{\resp}{respectively\ }
\newcommand{\Subsection}{\subsection}
\providecommand{\nobreakdash}{\nobreak\hbox{-}}
\setlength{\emergencystretch}{2em}
\multlinegap0pt
\usepackage{amssymb}
\usepackage{mathtools}
\newtheorem{proposition}{Proposition}
\title{Vector Invariance and Structural Closure of Julia-Type Iterations in Clifford Algebra}
\author{Orgest ZAKA}
\date{Source: arXiv:2605.20252v1 (CC BY 4.0)}
\begin{document}
\maketitle

\noindent\emph{Source and licence.} Vector Invariance and Structural Closure of Julia-Type Iterations in Clifford Algebra. Version arXiv:2605.20252v1 (CC BY 4.0), distributed by arXiv under the Creative Commons Attribution 4.0 International licence, http://creativecommons.org/licenses/by/4.0/, as recorded in the arXiv licence metadata for this exact version and shown on the abstract page https://arxiv.org/abs/2605.20252v1. Full licence notice: \texttt{licenses/arxiv-2605.20252-CC-BY-4.0.txt}.\par\smallskip

\noindent\textit{Editorial scope.} Counted unit: the article's proposition proposition (source0.tex lines 353--386). The article's complete original proof of it is delivered with it (source0.tex lines 387--668, one \texttt{proof} environment of 282 lines). No mathematics is written, repaired or re-derived: the statement and the proof are the article's own lines, in the article's own notation, and no retained line is substituted or re-flowed.  No further source-local context is needed: every cross-reference of every retained line resolves inside the two delivered spans.  Nothing was omitted from any retained span, so the package carries no omission record.  No retained line contains a citation, so the wrapper carries no bibliography and no bibliography entry.  No retained line contains a figure environment, an image-inclusion command or drawing code, so no image is delivered and the package declares no asset dependency.  Editorial formatting only, in addition: an AMS \texttt{amsart} preamble that restates the article's own declarations and macros used by the retained lines, namely the environments \texttt{proposition} on separate counters, together with the standard packages those lines need.  The retained environments print in this document as Proposition 1 in order of appearance, whereas the article prints the same items with the numbering of its own full document.  The complete native source of the article as distributed in the arXiv e-print for this exact version is bundled unchanged as \texttt{sources/200980/source0.tex}.
\begin{proposition}[Vector Invariance in $\mathbb{R}^{3}$]
Let
\[
\vec{x},\vec{n},\vec{c}\in \mathbb{R}^{3},
\]
where $\vec n$ is a unit vector satisfying
\[
\vec n\diamond \vec n =1.
\]

For every integer $p\ge 2$, define
\[
f(\vec{x})
=
(\vec{x}\diamond \vec n)^p\diamond \vec n+\vec c.
\]

Then
\[
f:\mathbb{R}^{3}\to \mathbb{R}^{3}
\]
is well-defined.

In particular,
\[
(\vec{x}\diamond \vec n)^p\diamond \vec n
\]
belongs to the grade-$1$ subspace of the Clifford algebra associated
with $\mathbb{R}^{3}$.

Therefore, despite the appearance of bivector and trivector terms in
intermediate computations, the final expression reduces to a vector in
$\mathbb{R}^{3}$.
\end{proposition}

\begin{proof}
Without loss of generality, we choose an orthonormal basis
$\{\vec e_1,\vec e_2,\vec e_3\}$ and take $\vec n=\vec e_1$.

Let
\[
\vec x = x_1\vec e_1 + x_2\vec e_2 + x_3\vec e_3.
\]

\textbf{Case $p=1$.}
We have
\[
(\vec x\diamond \vec e_1)\diamond \vec e_1
=
\vec x\diamond (\vec e_1\diamond \vec e_1)
=
\vec x,
\]
so the claim holds.

\medskip

\textbf{Case $p=2$.}
Thinking and acting in the same way in $\mathbb{R}^{3}$ we have
\begin{align*}
\left( \vec{x} \diamond \vec{e}_{1} \right)^{2} \diamond \vec{e}_{1}
&= \left( \vec{x} \diamond \vec{e}_{1} \right) \diamond 
   \left( \vec{x} \diamond \vec{e}_{1} \right) \diamond \vec{e}_{1} \\[4pt]
&= \left( \vec{x} \diamond \vec{e}_{1} \diamond \vec{x} \right)
   \diamond \left( \vec{e}_{1} \diamond \vec{e}_{1} \right) \\[4pt]
&= \left( \vec{x} \diamond \vec{e}_{1} \diamond \vec{x} \right)
   \left( \vec{e}_{1} \cdot \vec{e}_{1} + \vec{e}_{1} \wedge \vec{e}_{1} \right) \\[4pt]
&= \left( \vec{x} \diamond \vec{e}_{1} \diamond \vec{x} \right) .\\
& \text{(Since $\vec{e}_{1} \cdot \vec{e}_{1} = 1$ and $\vec{e}_{1} \wedge \vec{e}_{1} = 0$ )}\\[4pt]
%5\end{align*}
%\begin{align*}
%\text{Since } \vec{e}_{1} \cdot \vec{e}_{1} &= 1 
%\quad \text{and} \quad 
%\vec{e}_{1} \wedge \vec{e}_{1} = 0, \\[6pt]
&= \left( 
x_{1}\vec{e}_{1} + x_{2}\vec{e}_{2} + x_{3}\vec{e}_{3}
\right)
\diamond \vec{e}_{1}
\diamond 
\left(
x_{1}\vec{e}_{1} + x_{2}\vec{e}_{2} + x_{3}\vec{e}_{3}
\right) \\[6pt]
&= \Big[
x_{1}\left( \vec{e}_{1} \diamond \vec{e}_{1} \right)
+ x_{2}\left( \vec{e}_{2} \diamond \vec{e}_{1} \right)
+ x_{3}\left( \vec{e}_{3} \diamond \vec{e}_{1} \right)
\Big] \diamond
\left(
x_{1}\vec{e}_{1} + x_{2}\vec{e_2} + x_{3}\vec{e_3}
\right).\\[4pt]
& \text{(Since $\vec{e}_{1} \cdot \vec{e}_{1} = 1$ and $\vec{e_1} \wedge \vec{e_1} = 0$)}\\[4pt]
%\end{align*}
%\begin{align*}
&= \left(
x_{1}\vec{e}_{1} + x_{2}\vec{e}_{2} + x_{3}\vec{e}_{3}
\right)
\diamond \vec{e}_{1}
\diamond
\left(
x_{1}\vec{e}_{1} + x_{2}\vec{e}_{2} + x_{3}\vec{e}_{3}
\right) \\[6pt]
&= \Big[
x_{1}\left( \vec{e}_{1}\diamond \vec{e}_{1} \right)
+ x_{2}\left( \vec{e}_{2}\diamond \vec{e}_{1} \right)
+ x_{3}\left( \vec{e}_{3}\diamond \vec{e}_{1} \right)
\Big]
\diamond
\left(
x_{1}\vec{e}_{1} + x_{2}\vec{e}_{2} + x_{3}\vec{e}_{3}
\right) \\[6pt]
&= \Bigg\{ \Big[
x_{1}\left( \vec{e}_{1}\cdot \vec{e}_{1} + \vec{e}_{1}\wedge \vec{e}_{1} \right)
\Big] \\
&\quad + \Big[
x_{2}\left( \vec{e}_{2}\cdot \vec{e}_{1} + \vec{e}_{2}\wedge \vec{e}_{1} \right)
\Big] \\
&\quad + \Big[
x_{3}\left( \vec{e}_{3}\cdot \vec{e}_{1} + \vec{e}_{3}\wedge \vec{e}_{1} \right)
\Big]
\Bigg \}
\diamond
\left(
x_{1}\vec{e}_{1} + x_{2}\vec{e}_{2} + x_{3}\vec{e}_{3}
\right) \\[6pt]
&= \Big[
x_{1}
+ x_{2}\left( \vec{e}_{2}\wedge \vec{e}_{1} \right)
+ x_{3}\left( \vec{e}_{3}\wedge \vec{e}_{1} \right)
\Big]
\diamond
\left(
x_{1}\vec{e}_{1} + x_{2}\vec{e}_{2} + x_{3}\vec{e}_{3}
\right). \\[4pt]
&\text{(but $\vec{e}_{2}\cdot \vec{e}_{1}=\vec{e_{3}}\cdot \vec{e}_{1}=0$ and $\vec{e}_{1}\cdot \vec{e}_{1}=1$) }
\end{align*}
So,
\begin{align*}
\left(\vec{x}\diamond\vec{e_1}\right)^2\diamond\vec{e_1}&= x_{1}\Big[
x_{1}+x_{2}\left( \vec{e}_{2}\wedge \vec{e}_{1} \right)
+ x_{3}\left( \vec{e}_{3}\wedge \vec{e}_{1} \right)
\Big]\diamond \vec{e}_{1} \\[4pt]
&\quad + x_{2}\Big[
x_{1}+x_{2}\left( \vec{e}_{2}\wedge \vec{e}_{1} \right)
+ x_{3}\left( \vec{e}_{3}\wedge \vec{e}_{1} \right)
\Big]\diamond \vec{e}_{2} \\[4pt]
&\quad + x_{3}\Big[
x_{1}+x_{2}\left( \vec{e}_{2}\wedge \vec{e}_{1} \right)
+ x_{3}\left( \vec{e}_{3}\wedge \vec{e}_{1} \right)
\Big]\diamond \vec{e}_{3} \\[8pt]
&= \Big[
x_{1}^{2}
+ x_{1}x_{2}\left( \vec{e}_{2}\wedge \vec{e}_{1} \right)
+ x_{1}x_{3}\left( \vec{e}_{3}\wedge \vec{e}_{1} \right)
\Big]\diamond \vec{e}_{1} \\[4pt]
&\quad + \Big[
x_{1}x_{2}
+ x_{2}^{2}\left( \vec{e}_{2}\wedge \vec{e}_{1} \right)
+ x_{2}x_{3}\left( \vec{e}_{3}\wedge \vec{e}_{1} \right)
\Big]\diamond \vec{e}_{2} \\[4pt]
&\quad + \Big[
x_{1}x_{3}
+ x_{2}x_{3}\left( \vec{e}_{2}\wedge \vec{e}_{1} \right)
+ x_{3}^{2}\left( \vec{e}_{3}\wedge \vec{e}_{1} \right)
\Big]\diamond \vec{e}_{3} \\[4pt]
&= x_{1}^{2}\vec{e}_{1}
+ x_{1}x_{2}\left( \vec{e}_{2}\wedge \vec{e}_{1} \right)\diamond \vec{e}_{1}
+ x_{1}x_{3}\left( \vec{e}_{3}\wedge \vec{e}_{1} \right)\diamond \vec{e}_{1} \\[6pt]
&\quad + x_{1}x_{2}\vec{e}_{2}
- x_{2}^{2}\left( \vec{e}_{1}\wedge \vec{e}_{2} \right)\diamond \vec{e}_{2}
+ x_{2}x_{3}\left( \vec{e}_{3}\wedge \vec{e}_{1} \right)\diamond \vec{e}_{2} \\[6pt]
&\quad + x_{1}x_{3}\vec{e}_{3}
+ x_{2}x_{3}\left( \vec{e}_{2}\wedge \vec{e}_{1} \right)\diamond \vec{e}_{3}
+ x_{3}^{2}\left( \vec{e}_{3}\wedge \vec{e}_{1} \right)\diamond \vec{e}_{3}. \\[4pt]
&= \Big[
x_{1}^{2}\vec{e}_{1}
+ x_{1}x_{2}\left( \vec{e}_{2}\wedge \vec{e}_{1} \right)\diamond \vec{e}_{1}
+ x_{1}x_{3}\vec{e}_{3}
\Big] \\[6pt]
&\quad + \Big[
x_{1}x_{2}\vec{e}_{2}
- x_{2}^{2}\left( \vec{e}_{1}\wedge \vec{e}_{2} \right)\diamond \vec{e}_{2}
+ x_{2}x_{3}\left( \vec{e}_{3}\wedge \vec{e}_{1} \right)\diamond \vec{e}_{2}
\Big] \\[6pt]
&\quad + \Big[
x_{1}x_{3}\vec{e}_{3}
+ x_{2}x_{3}\left( \vec{e}_{2}\wedge \vec{e}_{1} \right)\diamond \vec{e}_{3}
- x_{3}^{2}\left( \vec{e}_{1}\wedge \vec{e}_{3} \right)\diamond \vec{e}_{3}
\Big].
\end{align*}

\textbf{Remark.} We use the identities:
\begin{align*}
\left( \vec{e}_{i} \wedge \vec{e}_{j} \right) \diamond \vec{e}_{j}
&= \vec{e}_{i}, \\[4pt]
\vec{e}_{i} \wedge \vec{e}_{j}
&= - \vec{e}_{j} \wedge \vec{e}_{i}, \\[4pt]
\left( \vec{e}_{i} \wedge \vec{e}_{j} \right) \diamond \vec{e}_{k}
&= \left( \vec{e}_{i} \wedge \vec{e}_{j} \right) \cdot \vec{e}_{k}
+ \left( \vec{e}_{i} \wedge \vec{e}_{j} \right) \wedge \vec{e}_{k}.
\end{align*}
And next we have:
\begin{align*}
(\vec{x}\diamond\vec{e_1})^2\diamond \vec{e_1}&= \Big[
x_{1}^{2}\vec{e}_{1}
+ x_{1}x_{2}\vec{e}_{2}
+ x_{1}x_{3}\vec{e}_{3}
\Big] \\[4pt]
&\quad + \Big[
x_{1}x_{2}\vec{e}_{2}
- x_{2}^{2}\vec{e}_{1}
+ x_{2}x_{3}\left( \vec{e}_{3}\wedge \vec{e}_{1} \right)\diamond \vec{e}_{2}
\Big] \\[4pt]
&\quad + \Big[
x_{1}x_{3}\vec{e}_{3}
+ x_{2}x_{3}\left( \vec{e}_{2}\wedge \vec{e}_{1} \right)\diamond \vec{e}_{3}
- x_{3}^{2}\vec{e}_{1}
\Big] \\[8pt]
&= x_{1}^{2}\vec{e}_{1}
+ x_{1}x_{2}\vec{e}_{2}
+ x_{1}x_{3}\vec{e}_{3}
+ x_{1}x_{2}\vec{e}_{2}
- x_{2}^{2}\vec{e}_{1} \\[4pt]
&\quad + x_{2}x_{3}\left( \vec{e}_{3}\wedge \vec{e}_{1} \right)\diamond \vec{e}_{2} \\[4pt]
&\quad + x_{1}x_{3}\vec{e}_{3}
+ x_{2}x_{3}\left( \vec{e}_{2}\wedge \vec{e}_{1} \right)\diamond \vec{e}_{3}
- x_{3}^{2}\vec{e}_{1} \\[4pt]
%\end{align*}
%\begin{align*}
&= \left( x_{1}^{2}-x_{2}^{2}-x_{3}^{2} \right)\vec{e}_{1}
+ \left( x_{1}x_{2}+x_{1}x_{2} \right)\vec{e}_{2}
+ \left( x_{1}x_{3}+x_{1}x_{3} \right)\vec{e}_{3} \\[4pt]
&\quad + x_{2}x_{3}
\Big[
\left( \vec{e}_{3}\wedge \vec{e}_{1} \right)\diamond \vec{e}_{2}
+ \left( \vec{e}_{2}\wedge \vec{e}_{1} \right)\diamond \vec{e}_{3}
\Big] \\[8pt]
&= \left( x_{1}^{2}-x_{2}^{2}-x_{3}^{2} \right)\vec{e}_{1}
+ 2x_{1}x_{2}\vec{e}_{2}
+ 2x_{1}x_{3}\vec{e}_{3} \\[6pt]
&\quad + x_{2}x_{3}
\Big[
\left( \vec{e}_{2}\cdot \vec{e}_{3} \right)\vec{e}_{1}
- \left( \vec{e}_{2}\cdot \vec{e}_{1} \right)\vec{e}_{3}
+ \vec{e}_{3}\wedge \vec{e}_{1}\wedge \vec{e}_{2} \\[4pt]
&\qquad + \left( \vec{e}_{3}\cdot \vec{e}_{2} \right)\vec{e}_{1}
- \left( \vec{e}_{3}\cdot \vec{e}_{1} \right)\vec{e}_{2}
+ \vec{e}_{2}\wedge \vec{e}_{1}\wedge \vec{e}_{3}
\Big] \\[8pt]
&= \left( x_{1}^{2}-x_{2}^{2}-x_{3}^{2} \right)\vec{e}_{1}
+ 2x_{1}x_{2}\vec{e}_{2}
+ 2x_{1}x_{3}\vec{e}_{3} \\[6pt]
&\quad + x_{2}x_{3}
\Big[
\vec{e}_{3}\wedge \vec{e}_{1}\wedge \vec{e}_{2}
+ \vec{e}_{2}\wedge \vec{e}_{1}\wedge \vec{e}_{3}
\Big].
\end{align*}
But we have that
\begin{align*}
\vec{e}_{3} \wedge \vec{e}_{1} \wedge \vec{e}_{2}
&= - \vec{e}_{2} \wedge \vec{e}_{1} \wedge \vec{e}_{3}, \\[6pt]
&\text{and hence the expression reduces to } \\[6pt]
&= \left( x_{1}^{2}-x_{2}^{2}-x_{3}^{2} \right)\vec{e}_{1}
+ 2x_{1}x_{2}\vec{e}_{2}
+ 2x_{1}x_{3}\vec{e}_{3}.
\end{align*}
which is a vector $\mathbb{R}^3$.


\textbf{Inductive step.}
Assume that for some $p\ge 2$,
\[
(\vec x\diamond \vec e_1)^{p-1}\diamond \vec e_1 \in \mathbb{R}^3.
\]

Since any vector in $\mathbb{R}^3$ can be written as
\[
\vec v = v_1\vec e_1 + v_2\vec e_2 + v_3\vec e_3,
\]
we write
\[
(\vec x\diamond \vec e_1)^{p-1}\diamond \vec e_1
=
a_1\vec e_1 + a_2\vec e_2 + a_3\vec e_3.
\]

Then
\begin{align*}
(\vec x\diamond \vec e_1)^p\diamond \vec e_1
&=
(\vec x\diamond \vec e_1)\diamond (a_1\vec e_1 + a_2\vec e_2 + a_3\vec e_3).
\end{align*}

Now use
\[
\vec x\diamond \vec e_1
=
x_1 + x_2(\vec e_2\wedge \vec e_1) + x_3(\vec e_3\wedge \vec e_1).
\]

Expanding term by term, we get:
\begin{itemize}
\item  scalar $\cdot$ vector gives vector,
\item  bivector $\diamond \vec e_1$ gives vector,
\item  and all mixed products reduce via identities
\[
(\vec e_i\wedge \vec e_1)\diamond \vec e_1 = \vec e_i.
\]
\end{itemize}
Hence every term in the expansion is a linear combination of
$\vec e_1,\vec e_2,\vec e_3$, so:
\[
(\vec x\diamond \vec e_1)^p\diamond \vec e_1 \in \mathbb{R}^3.
\]

Therefore, by induction, the result holds for all $p\ge 1$.
\end{proof}

\end{document}
